\chapter{Graph}

\section{Fundamentals}
	% \kactlimport{BellmanFord.h}
	% \kactlimport{FloydWarshall.h}
	% \kactlimport{TopoSort.h}

\section{Network flow}
	\kactlimport{PushRelabel.h}
	\kactlimport{MinCostMaxFlow.h}
	\kactlimport{EdmondsKarp.h}
	\kactlimport{Dinic.h}
		
	\subsection*{Circulation and Lower Bounds}
		\textbf{Circulation:}
		Each vertex $v$ has demand
		$f^{in}(v)-f^{out}(v)=d_v$.
		Add super source $SS$, super sink $TT$:
		\begin{itemize}
		\setlength{\itemsep}{0pt}
		\item $d_v<0$: add $SS\to v$, cap $-d_v$.
		\item $d_v>0$: add $v\to TT$, cap $d_v$.
		\end{itemize}
		Feasible iff $\sum_v d_v=0$ and
		$\operatorname{maxflow}(SS,TT)=\sum_{d_v>0}d_v$.

		\textbf{Lower bounds:}
		For each edge $(u,v)$ with bounds $[l,c]$,
		take $l$ units of flow first:
		\[
		c'=c-l,\quad d_u\mathrel{+}=l,\quad d_v\mathrel{-}=l.
		\]
		Solve circulation with updated demands.
		Recover original edge flow: $f=l+f'$.

		\textbf{Max flow with lower bounds:}
		Initially $d_v=0$ for all vertices.
		After removing lower bounds, add $t\to s$
		with capacity $\infty$ and solve circulation.
		Let $F_0=f(t,s)$.
		Remove $SS,TT$ and $t\to s$, keeping
		the residual graph of original edges.
		\[
		F_{\max}=F_0+\operatorname{maxflow}(s,t).
		\]
		Original edge flow: $f=l+f'$ (final flow).
		
		\textbf{Time:} Two max-flow computations.
	
	\kactlimport{MinCut.h}
	\kactlimport{GlobalMinCut.h}
	\kactlimport{GomoryHu.h}

\section{Matching}
	\kactlimport{HopcroftKarp.h}
	\kactlimport{DFSMatching.h}
	\kactlimport{MinimumVertexCover.h}
	\kactlimport{WeightedMatching.h}
	% \kactlimport{GeneralMatching.h}
	\kactlimport{Blossom.h}

\section{DFS algorithms}
	% \kactlimport{SCC.h}
	\kactlimport{BiconnectedComponents.h}
	\kactlimport{2sat.h}
	% \kactlimport{EulerWalk.h}

\section{Coloring}
	\kactlimport{EdgeColoring.h}

\section{Heuristics}
	\kactlimport{MaximalCliques.h}
	\kactlimport{MaximumClique.h}
	\kactlimport{MaximumIndependentSet.h}

\section{Trees}
	% \kactlimport{BinaryLifting.h}
	\kactlimport{LCA.h}
	\kactlimport{CompressTree.h}
	% \kactlimport{HLD.h}
	\kactlimport{LinkCutTree.h}
	\kactlimport{DirectedMST.h}

\section{Math}
	\subsection{Number of Spanning Trees}
		% I.e. matrix-tree theorem.
		% Source: https://en.wikipedia.org/wiki/Kirchhoff%27s_theorem
		% Test: stress-tests/graph/matrix-tree.cpp
		Create an $N\times N$ matrix \texttt{mat}, and for each edge $a \rightarrow b \in G$, do
		\texttt{mat[a][b]--, mat[b][b]++} (and \texttt{mat[b][a]--, mat[a][a]++} if $G$ is undirected).
		Remove the $i$th row and column and take the determinant; this yields the number of directed spanning trees rooted at $i$
		(if $G$ is undirected, remove any row/column).
	\par\vspace{12pt}
	\subsection{Erdős–Gallai theorem}
		% Source: https://en.wikipedia.org/wiki/Erd%C5%91s%E2%80%93Gallai_theorem
		% Test: stress-tests/graph/erdos-gallai.cpp
		A simple graph with node degrees $d_1 \ge \dots \ge d_n$ exists iff $d_1 + \dots + d_n$ is even and for every $k = 1\dots n$,
		\[ \sum _{i=1}^{k}d_{i}\leq k(k-1)+\sum _{i=k+1}^{n}\min(d_{i},k). \]
